\documentclass[10pt,a4paper]{amsart}
\usepackage[margin=19mm]{geometry}
\usepackage{amsmath,amssymb,mathtools}
\usepackage[scr=rsfs,cal=euler]{mathalfa}
\usepackage{xfrac}
\usepackage[unicode,hidelinks,bookmarks=false]{hyperref}
\newcommand{\ie}{i.e.\ }
\newcommand{\cf}{cf.\ }
\newcommand{\resp}{respectively\ }
\newcommand{\Subsection}{\subsection}
\providecommand{\nobreakdash}{\nobreak\hbox{-}}
\setlength{\emergencystretch}{2em}
\multlinegap0pt
\usepackage{amssymb}
\usepackage{mathtools}
\usepackage{tikz}
\usepackage{caption}
\newtheorem{lemma}{Lemma}
\DeclareMathOperator{\per}{per}
\title{A Permanental Analog of the Rank--Nullity Theorem for Symmetric Matrices}
\author{Priyanshu Pant, Surabhi Chakrabartty, Ranveer Singh}
\date{Source: arXiv:2507.01344v4 (CC BY 4.0)}
\begin{document}
\maketitle

\noindent\emph{Source and licence.} A Permanental Analog of the Rank--Nullity Theorem for Symmetric Matrices. Version arXiv:2507.01344v4 (CC BY 4.0), distributed by arXiv under the Creative Commons Attribution 4.0 International licence, http://creativecommons.org/licenses/by/4.0/, as recorded in the arXiv licence metadata for this exact version and shown on the abstract page https://arxiv.org/abs/2507.01344v4. Full licence notice: \texttt{licenses/arxiv-2507.01344-CC-BY-4.0.txt}.\par\smallskip

\noindent\textit{Editorial scope.} Counted unit: the article's lemma lemma:app-principal-existence (source0.tex lines 811--816). The article's complete original proof of it is delivered with it (source0.tex lines 818--881, one \texttt{proof} environment of 64 lines). No mathematics is written, repaired or re-derived: the statement and the proof are the article's own lines, in the article's own notation, and no retained line is substituted or re-flowed.  Retained uncounted as the source-local context the proof needs: lines 778--782.  Nothing was omitted from any retained span, so the package carries no omission record.  No retained line contains a citation, so the wrapper carries no bibliography and no bibliography entry.  No retained line contains a figure environment, an image-inclusion command or drawing code, so no image is delivered and the package declares no asset dependency.  Editorial formatting only, in addition: an AMS \texttt{amsart} preamble that restates the article's own declarations and macros used by the retained lines, namely the environments \texttt{lemma} on separate counters, together with the standard packages those lines need.  The retained environments print in this document as Lemma 1 in order of appearance, whereas the article prints the same items with the numbering of its own full document.  The complete native source of the article as distributed in the arXiv e-print for this exact version is bundled unchanged as \texttt{sources/180870/source0.tex}.
\begin{lemma}\label{lemma:app-bijection-constraint}
  Let \( A \) be a non-negative symmetric matrix with \( \rho_{\per}(A) = k \). 
  Then, any permutation $\sigma$ corresponding to which a submatrix $A[I, J]$ of order $k$ has non-zero permanent, must satisfy the following condition:
    \[ \sigma(i) \notin I \implies i \in J\]
\end{lemma}

\begin{lemma}\label{lemma:app-principal-existence}
Let \( A \) be a non-negative symmetric matrix of order \( n \) with \( \rho_{\per}(A) = k \). Then, there exists a principal submatrix \( A[S, S] \) of order \( k \) such that
\[
\per(A[S, S]) \neq 0.
\]
\end{lemma}

\begin{proof}

Since $\rho_{\per}(A)=k$, so there exists a submatrix \( A[I, J] \) (with \( I, J \subseteq \{1, 2, \dots, n\} \) and \( |I| = |J| = k \)) such that
\(
\per(A[I, J]) \neq 0.
\)

If $I=J$, then the proof is complete. So, we will consider the case when $I \neq J$. Now, as $\per(A[I, J]) \neq 0$, by definition of permanent there exists a bijection $\sigma \in S_n$ such that for each $i \in I$, there exists a $ \sigma(i) \in J$ and 
$ A(i,\sigma(i)) \neq 0. $

Now, we show that the principal submatrix \( A[I, I] \) has a nonzero permanent. Define a permutation \( \pi: I\to I \) as follows.
\[
\pi(i) = 
\begin{cases} 
\sigma(i), & \text{if } \sigma(i) \in I, \\[1mm]
\sigma^{-1}(i), & \text{if } \sigma(i) \notin I.
\end{cases}
\]
Since, by Lemma \ref{lemma:app-bijection-constraint} if \(\sigma(i) \notin I \implies i \in J \), \( \pi(i) \) is well-defined for all \( i \in I\). \\
Now we verify two claims about \( \pi \):


\textbf{Claim 1: }  \( \pi \) is a bijection on \( I \).


We consider three cases for \( i_1, i_2 \in I \):

\begin{enumerate}
    \item  If \( \sigma(i_1), \sigma(i_2) \in I \), then \( \pi(i_1) = \sigma(i_1) \) and \( \pi(i_2) = \sigma(i_2) \). Since \( \sigma \) is injective, \( \pi(i_1) = \pi(i_2) \) implies \( i_1 = i_2 \).
    
    \item  If \( \sigma(i_1), \sigma(i_2) \notin I \), then \( \pi(i_1) = \sigma^{-1}(i_1) \) and \( \pi(i_2) = \sigma^{-1}(i_2) \). Since \( \sigma^{-1} \) is injective, again \( \pi(i_1) = \pi(i_2) \) implies \( i_1 = i_2 \).
    
    \item If \( \sigma(i_1) \in I \) but \( \sigma(i_2) \notin I \), then \( \pi(i_1) = \sigma(i_1) \) and \( \pi(i_2) = \sigma^{-1}(i_2) \). Suppose \( \pi(i_1) = \pi(i_2) = c \). Then \( \sigma(i_1) = c \) and \( \sigma^{-1}(i_2) = c \), so \( \sigma(c) = i_2 \). By the defining property of \( \sigma \), we have \( A(i_1, \sigma(i_1)) = A(i_1, c) \neq 0 \). Since \( A \) is symmetric, \( A(i_2, c) = A(c, i_2) = A(i_2, \sigma(i_1)) \neq 0 \). This contradicts the bijectivity of \( \sigma \) unless \( i_1 = i_2 \).
\end{enumerate}
Thus, \( \pi \) is injective (and hence bijective on \( I \)).


\textbf{Claim 2:} For each \( i \in I \), \( A(i,\pi(i)) \neq 0 .\)  

There are two cases to consider:
\begin{enumerate}
    \item If \( \sigma(i) \in I \) then \( \pi(i) = \sigma(i) \), and by definition of \( \sigma \), we have
    \(
    A(i, \pi(i)) = A(i, \sigma(i)) \neq 0.
    \)
    
    \item If \( \sigma(i) \notin I \), then by Lemma \ref{lemma:app-bijection-constraint}, it follows that \( i \in J \). Thus, there exists a unique \( i' \in I \) such that \( \sigma(i') = i \) that is, \( i' = \sigma^{-1}(i) \). By the definition of \( \sigma \),
    \(
    A(i', \sigma(i)) = A(i', i) \neq 0.
    \)
    Since \( A \) is symmetric, we also have \( A(i, i') = A(i, \sigma^{-1}(i)) \neq 0 \). In this case, \( \pi(i) = \sigma^{-1}(i) \), so once again \( A(i, \pi(i)) \neq 0 \).
\end{enumerate}


Thus, for every \( i \in I \), \( A(i,\pi(i))\neq 0 \). It follows that
\(
\prod_{i \in I} A(i,\pi(i)) \neq 0,
\)
and hence,
\(
\per(A[I, I]) \neq 0.
\)
This shows that \( A[I,I] \) is a principal submatrix of order \( k \) with nonzero permanent, which completes the proof. \\
\end{proof}

\end{document}
